\chapter{Weak Interactions, Virial Corrections, and Inexact Differentials}

\lectureoverview{We begin with the operational meaning of measurement and
temperature, then ask how a dilute gas departs from the ideal-gas law when
short-range interactions are weakly turned on.  The partition function gives
the first interaction corrections to energy and pressure and identifies their
range of validity.  A final mathematical interlude distinguishes state
functions from path-dependent heat and work by means of exact differentials,
closed loops, and a concrete ideal-gas test.}

\section{What an Instrument Actually Measures}

No laboratory device reaches directly into nature and reports an abstract
mathematical variable.  A spring balance reports an extension; a thermometer
reports the height of a liquid column.  We call the first a force measurement
and the second a temperature measurement only because a reproducible relation
has been established between the visible response and the desired quantity.

That relation has a limited range.  A spring obeys Hooke's law only while it
remains elastic, and a sufficiently large force eventually deforms or breaks
it.  More generally, if a smooth readout \(R\) depends on a physical variable
\(X\), then near a chosen operating point
\begin{equation}
 \Delta R \simeq
 \left.\frac{dR}{dX}\right|_{X_0}\Delta X.
 \label{eq:sm6-local-linear-response}
\end{equation}
Small changes in the readout are proportional to small changes in the target
variable even when the full response curve is nonlinear.  This local
linearization is the ordinary logic of calibration.

The readout can also depend on nuisance variables.  A spring may respond
weakly to temperature or pressure.  The height of mercury in a thermometer
could respond to orientation through gravity, to properties of the glass, or
to slow changes in the materials.  A measurement protocol does not prove
that these dependencies are absent; it holds them fixed or shows that they
are negligible on the required scale.

\begin{classroomqa}
\audiencequestion{Could the local gravitational field, or the direction in
which the thermometer is held, change the reading?}
\lecturerresponse{Yes in principle.  If orientation has an appreciable
effect, the instruction to hold the thermometer vertically becomes part of
the instrument's definition.  Near Earth's surface the gravitational field
is sufficiently stable that the remaining variation is normally negligible.
\lecturetimestamp{00:04:08--00:05:00}}
\end{classroomqa}

\begin{classroomqa}
\audiencequestion{Could aging or chemical change slowly alter the
calibration?}
\lecturerresponse{It could.  Mercury and glass are sufficiently inert for
ordinary measurements, while changes on fantastically long time scales are
irrelevant to the experiment.  Calibration is always a statement about a
specified instrument, range, protocol, and duration.
\lecturetimestamp{00:05:03--00:06:11}}
\end{classroomqa}

For a Celsius thermometer, two reproducible reference states supply a
convenient scale.  If \(h_f\) and \(h_b\) are the column heights at the
freezing and boiling points of water, respectively, then within a range where
the response is nearly linear,
\begin{equation}
 T_{\mathrm C}=100\,\frac{h-h_f}{h_b-h_f}.
 \label{eq:sm6-celsius-calibration}
\end{equation}
The choice \(T_{\mathrm C}=0\) at the freezing point is conventional.  It is
not the state of zero thermal energy, just as zero degrees Fahrenheit on a
winter day is not absolute zero.

\begin{figure}[t]
\centering
\includegraphics[width=0.90\textwidth]{lecture_06_figure_03.png}
\caption{A smooth but nonlinear thermometer response, with a short interval
marked out for linear calibration.}
\lectureframe{00:07:55}
\end{figure}

\section{Equilibrium and Temperature as an Energy Scale}

A thermometer measures the temperature of its surroundings only after the
thermometer and the nearby medium have reached thermal equilibrium.  Before
then, its reading is merely changing toward equilibrium.  The practical test
is elementary: wait until the reading stops changing.  The simplicity of
that rule hides the physical processes that set the waiting time.

\begin{classroomqa}
\audiencequestion{Does equilibration involve only a boundary layer around
the bulb, and how long should one wait?}
\lecturerresponse{Energy must cross the glass and be exchanged with the
surrounding medium, so the time depends on the thermometer, its thermal
conductivity, and the medium.  The immediate surroundings must also be
representative of a uniform equilibrium state.  Experimentally, the reading
is accepted when it has ceased to drift.
\lecturetimestamp{00:08:45--00:10:13}}
\end{classroomqa}

Temperature and density play different roles in this waiting time.  A very
dilute gas can be very hot: the few molecules present may carry large kinetic
energies.  Yet those molecules strike the thermometer rarely, so equilibration
can be slow.  Lowering density by a factor of one hundred or one thousand
while holding temperature fixed does not lower the molecular energy scale;
it lowers the rate at which that energy is communicated to the instrument.

\begin{classroomqa}
\audiencequestion{Can the radiation in a room have a different temperature
from the air?}
\lecturerresponse{Yes.  Matter and radiation share one temperature only when
they are in thermal equilibrium with one another.  Neutral air couples
weakly to much of the optical radiation in an illuminated room, and the
photon population need not be an equilibrium blackbody field at the air
temperature.
\lecturetimestamp{00:11:26--00:12:20}}
\end{classroomqa}

Once equilibrium is established, temperature is best regarded as an energy
scale.  We continue to use units in which \(k_B=1\), so \(T\) itself has
dimensions of energy.  For a monatomic ideal gas,
\begin{equation}
 E_0=\frac{3}{2}NT.
 \label{eq:sm6-ideal-energy}
\end{equation}
The thermal energy per molecule is fixed by \(T\); the corresponding speed
still depends on the molecular mass.

\section{Turning On a Weak Interaction}

\begin{classroomqa}
\audiencequestion{Under what conditions does the ideal-gas approximation
break down, and can the breakdown be shown mathematically?}
\lecturerresponse{Start from the solved noninteracting gas and turn on a
small interaction.  Expand in that change, calculate the first correction,
and compare it with the ideal term.  The approximation fails when the
correction is no longer small.
\lecturetimestamp{00:12:25--00:14:40}}
\end{classroomqa}

This is the standard perturbative strategy.  If a problem is soluble at
\(\lambda=0\), introduce a coupling \(\lambda\), write
\begin{equation}
 F(\lambda)=F(0)+\lambda F_1+\lambda^2F_2+\cdots,
 \label{eq:sm6-perturbative-template}
\end{equation}
and ask whether each retained term is smaller than the preceding one.  Here
the reference problem is the ideal gas.  The physical smallness comes from a
dilute sample, a short interaction range compared with the mean separation,
and a potential weak enough for the first interaction term to be useful.
High density or a long-range force can invalidate the approximation even
when any one encounter seems mild.

Let \(x_i\in\mathbb R^3\) and \(p_i\in\mathbb R^3\) be the position and
momentum of molecule \(i\).  For identical molecules of mass \(m\), take
\begin{equation}
 H(X,P)=
 \sum_{i=1}^{N}\frac{p_i^2}{2m}
 +\sum_{i<j}u\!\left(|x_i-x_j|\right),
 \label{eq:sm6-interacting-hamiltonian}
\end{equation}
where \(X=(x_1,\ldots,x_N)\) and \(P=(p_1,\ldots,p_N)\).  The restriction
\(i<j\) excludes self-pairs and counts \((i,j)\) only once:
\begin{equation}
 \sum_{i<j}1=\binom{N}{2}=\frac{N(N-1)}{2}.
 \label{eq:sm6-pair-count}
\end{equation}

\begin{classroomqa}
\audiencequestion{Why not sum independently over both particle labels?}
\lecturerresponse{The pair \((3,7)\) is the same physical pair as
\((7,3)\), and a particle is not paired with itself.  Restricting the sum to
\(i<j\) enforces both facts.  The mass symbol is kept as \(m\), while
\(i,j\) label particles, to avoid the notational collision that can occur
when \(m\) is used for both purposes.
\lecturetimestamp{00:17:52--00:19:12}}
\end{classroomqa}

Write the total configurational potential as
\begin{equation}
 U(X)=\sum_{i<j}u\!\left(|x_i-x_j|\right).
 \label{eq:sm6-total-potential}
\end{equation}
The interaction is assumed translationally invariant and short-ranged:
\(u(r)\) becomes negligible when \(r\) is large compared with a molecular
length.

\section{The Integrated Pair Strength}

Only one integral of the weak pair potential enters the first correction:
\begin{equation}
 U_0\equiv\int_{\mathbb R^3}d^3r\,u(r).
 \label{eq:sm6-u0-definition}
\end{equation}
Its units are energy times volume.  It combines the typical strength of the
interaction with the volume over which that interaction is appreciable.

\begin{figure}[t]
\centering
\includegraphics[width=0.90\textwidth]{lecture_06_figure_04.png}
\caption{The short-range pair potential compressed into the integrated
strength \(U_0=\int d^3r\,u(r)\).}
\lectureframe{00:25:15}
\end{figure}

To see why \(U_0\) is useful, consider one pair:
\begin{equation}
 I_2=\int_V d^3x_1\,d^3x_2\,
 u\!\left(|x_1-x_2|\right).
 \label{eq:sm6-two-particle-integral}
\end{equation}
Introduce a relative coordinate \(r=x_2-x_1\).  Away from the boundary, the
relative integral gives \(U_0\), while the common position of the pair ranges
through the full container.  Thus
\begin{equation}
 I_2=V U_0+\text{surface corrections}.
 \label{eq:sm6-pair-integral-vu0}
\end{equation}
The correction comes from particles within one interaction range of the
wall, where part of the relative-coordinate neighborhood lies outside the
box.  Its fractional importance scales as a surface-to-volume ratio and
vanishes in the macroscopic limit.

\begin{classroomqa}
\audiencequestion{Does \(U_0\) refer to one selected pair, and is one
particle kept at a fixed separation from the other?}
\lecturerresponse{It refers to the common interaction law for any one pair.
Hold the first particle fixed and integrate the second over every relative
position.  The separation is not fixed during this integral.  Only the
nearby interaction region contributes appreciably, producing \(U_0\).
\lecturetimestamp{00:26:03--00:27:33}}
\end{classroomqa}

\begin{classroomqa}
\audiencequestion{Why does the second position integral give the full
container volume while the first does not?}
\lecturerresponse{After the relative configuration has been integrated into
\(U_0\), the whole pair may be translated through the sample.  That common
translation supplies the large factor \(V\).  The interaction neighborhood
is a molecular-scale volume already contained inside \(U_0\), not another
factor equal to the container volume.
\lecturetimestamp{00:27:48--00:29:21}}
\end{classroomqa}

\begin{classroomqa}
\audiencequestion{Could the value depend on which point is held fixed,
especially near a wall?}
\lecturerresponse{Translational invariance makes the bulk value independent
of the fixed point.  Near a wall the integral is slightly different because
some neighboring positions are unavailable.  Only a surface layer of
molecular thickness is affected, so the error is negligible when the box is
macroscopic.
\lecturetimestamp{00:29:24--00:31:22}}
\end{classroomqa}

\begin{classroomqa}
\audiencequestion{Is the result the potential energy integrated over every
two-particle configuration, and why is it proportional to only one large
volume?}
\lecturerresponse{Yes.  Its units are energy times volume squared, since
\(V U_0\) contains one container volume and one molecular interaction
volume.  Only the first of those grows with the size of the sample.
\lecturetimestamp{00:31:31--00:32:04}}
\end{classroomqa}

\begin{classroomqa}
\audiencequestion{Does choosing a sphere, which minimizes area at fixed
volume, change the bulk answer?}
\lecturerresponse{It changes only the small boundary correction.  For any
ordinary macroscopic shape, the number of particles within a molecular
distance of the surface is negligible compared with the number in the bulk.
The surface becomes important only when the container itself approaches
molecular scales.
\lecturetimestamp{00:32:04--00:33:58}}
\end{classroomqa}

\begin{classroomqa}
\audiencequestion{Have collisions been omitted by replacing particle
trajectories with a potential?}
\lecturerresponse{No.  The force is the derivative of the potential and is
precisely what produces the collisions.  Equilibrium statistical mechanics
does not require us to follow those collisions one by one; their effect is
encoded in the energy and therefore in the Boltzmann weight.
\lecturetimestamp{00:33:59--00:35:29}}
\end{classroomqa}

\begin{classroomqa}
\audiencequestion{What is lost by retaining pair forces but not genuine
three-body interactions?}
\lecturerresponse{A microscopic energy can contain genuine three-body
terms, and higher orders of a pair potential also generate many-particle
correlations.  In a dilute gas, however, three specified particles are much
less likely than two to occupy one molecular neighborhood.  Pair effects
therefore give the first density correction; the higher clusters enter at
higher powers of density.
\lecturetimestamp{00:35:29--00:38:14}}
\end{classroomqa}

\section{The Partition Function on Autopilot}

The advantage of the statistical method is now visible.  Once the
Hamiltonian is specified, there is no need to guess the detailed velocity
distribution near walls or during collisions.  We write the canonical
partition function and follow its algebra:
\begin{equation}
 Z=
 \frac{1}{N!h^{3N}}
 \int d^{3N}P\,d^{3N}X\,
 e^{-\beta K(P)}e^{-\beta U(X)},
 \qquad
 K(P)=\sum_i\frac{p_i^2}{2m}.
 \label{eq:sm6-full-partition-function}
\end{equation}
The factor \(h^{3N}\) makes the classical phase-space measure dimensionless,
and \(N!\) removes the permutation overcounting of identical particles.

\begin{figure}[t]
\centering
\includegraphics[width=0.90\textwidth]{lecture_06_figure_05.png}
\caption{The many-particle energy and the corresponding kinetic and
configurational factors in the partition function.}
\lectureframe{00:40:35}
\end{figure}

Momentum and position separate.  Multiplying and dividing the
configurational integral by \(V^N\) isolates the ideal-gas result:
\begin{align}
 Z
 &=
 \left[
 \frac{V^N}{N!h^{3N}}
 \int d^{3N}P\,e^{-\beta K(P)}
 \right]
 \left[
 \frac{1}{V^N}
 \int d^{3N}X\,e^{-\beta U(X)}
 \right] \notag\\
 &\equiv Z_0\,C(\beta,V,N).
 \label{eq:sm6-z-factorization}
\end{align}
The familiar Gaussian factor is
\begin{equation}
 \int d^{3N}P\,e^{-\beta K(P)}
 =\left(\frac{2\pi m}{\beta}\right)^{3N/2}.
 \label{eq:sm6-gaussian-factor}
\end{equation}

\begin{figure}[t]
\centering
\includegraphics[width=0.90\textwidth]{lecture_06_figure_06.png}
\caption{The ideal-gas partition function isolated from the normalized
configurational factor.}
\lectureframe{00:43:45}
\end{figure}

The \(N!\) factor does not affect the normalized ratio \(C\), nor the
particular fixed-\(N\) derivatives used below, but it is not dispensable from
the thermodynamics as a whole: it is essential for correct state counting
and extensive entropy.

To first order in the interaction,
\begin{equation}
 e^{-\beta U(X)}
 =1-\beta U(X)+O(U^2).
 \label{eq:sm6-boltzmann-expansion}
\end{equation}
Consequently,
\begin{align}
 C
 &=
 1-\frac{\beta}{V^N}
 \sum_{i<j}
 \int d^{3N}X\,u(|x_i-x_j|)
 +O(U^2) \notag\\
 &=
 1-\frac{\beta}{V^N}
 \frac{N(N-1)}{2}
 V^{N-2}(V U_0)+O(U^2) \notag\\
 &=
 1-\beta\frac{N(N-1)}{2V}U_0+O(U^2).
 \label{eq:sm6-configurational-correction}
\end{align}
Each pair has the same \emph{integral}, not the same instantaneous
separation or potential energy.  The remaining \(N-2\) position integrals
each contribute one factor of \(V\).

\begin{figure}[t]
\centering
\includegraphics[width=0.90\textwidth]{lecture_06_figure_07.png}
\caption{Pair counting, spectator-coordinate volumes, and the first
configurational correction \(1-\beta N^2U_0/(2V)\).}
\lectureframe{00:53:40}
\end{figure}

\begin{classroomqa}
\audiencequestion{Does the leading \(1\) assume that the potential vanishes
at zero separation?}
\lecturerresponse{No.  It is the zeroth-order term in the exponential
series, independent of the value of \(u(0)\).  The potential enters in the
next term, \(-\beta U(X)\).
\lecturetimestamp{00:46:19--00:46:48}}
\end{classroomqa}

\begin{classroomqa}
\audiencequestion{Why may every pair be replaced by the same expression if
different pairs have different separations?}
\lecturerresponse{Their pointwise values are not being equated.  Every term
is integrated over the same full configuration space, and particle labels
do not distinguish one pair from another.  The integrals have identical
form and identical value.
\lecturetimestamp{00:50:00--00:51:07}}
\end{classroomqa}

\begin{classroomqa}
\audiencequestion{Is the number of pairs \(N(N-1)/2\) or
\(N(N+1)/2\)?}
\lecturerresponse{It is \(N(N-1)/2\): after selecting one particle there are
only \(N-1\) distinct partners, and division by two removes the reversed
ordering.  For macroscopic \(N\), replacing \(N(N-1)\) by \(N^2\) changes
only a relative \(1/N\) correction.
\lecturetimestamp{00:48:33--00:49:56}}
\end{classroomqa}

Since \(\log(1-x)=-x-x^2/2-x^3/3-\cdots\), not an alternating series, the
first-order logarithm is
\begin{equation}
 \log Z
 =
 \log Z_0
 -\beta\frac{N(N-1)}{2V}U_0
 +O(U^2).
 \label{eq:sm6-logz-correction}
\end{equation}
At large \(N\) we may replace \(N(N-1)\) by \(N^2\) where no finite-\(N\)
distinction is needed.

\begin{remark}
The controlled low-density formulation uses the Mayer function
\(f(r)=e^{-\beta u(r)}-1\) and the second virial coefficient
\begin{equation}
 B_2(T)=-\frac{1}{2}\int d^3r\,f(r).
 \label{eq:sm6-second-virial}
\end{equation}
For a weak potential, \(f(r)\simeq-\beta u(r)\), so
\(B_2\simeq\beta U_0/2\) and the result below is recovered.  This form is
needed for a hard repulsive core, where expanding \(e^{-\beta u}\) directly
is not valid.
\end{remark}

\section{Energy, Pressure, and the Meaning of the Correction}

The logarithm of \(Z\) is useful because its derivatives are measurable.
Assuming the microscopic potential is independent of temperature,
\begin{align}
 E
 &=-\left(\frac{\partial\log Z}{\partial\beta}\right)_{V,N} \notag\\
 &=
 \frac{3}{2}NT
 +\frac{N(N-1)}{2V}U_0+O(U^2) \notag\\
 &\simeq
 N\left(\frac{3}{2}T+\frac{1}{2}\rho U_0\right),
 \qquad
 \rho=\frac{N}{V}.
 \label{eq:sm6-energy-correction}
\end{align}
The apparently nonextensive \(N^2\) is accompanied by \(1/V\).  At fixed
density, \(N^2/V=N\rho\), so the total energy remains proportional to the
size of the system.  The extra energy per particle is proportional to
density because a chosen molecule acquires pair energy only when another
molecule lies nearby.

\begin{figure}[t]
\centering
\includegraphics[width=0.90\textwidth]{lecture_06_figure_08.png}
\caption{The first interaction correction reorganized as an extensive
energy \(N[3T/2+\rho U_0/2]\).}
\lectureframe{01:03:25}
\end{figure}

Pressure is easier to measure.  Using the result of the preceding chapter,
\begin{equation}
 P=T\left(\frac{\partial\log Z}{\partial V}\right)_{T,N},
 \label{eq:sm6-pressure-from-z}
\end{equation}
we obtain
\begin{align}
 P
 &=
 \frac{NT}{V}
 +T\frac{\partial}{\partial V}
 \left(-\beta\frac{N(N-1)}{2V}U_0\right)
 +O(U^2) \notag\\
 &=
 \rho T+\frac{1}{2}\rho^2U_0
 +O\!\left(\frac{\rho U_0}{N},U^2\right).
 \label{eq:sm6-pressure-correction}
\end{align}

\begin{figure}[t]
\centering
\includegraphics[width=0.90\textwidth]{lecture_06_figure_09.png}
\caption{Pressure from the Helmholtz derivative, including the first
quadratic-density correction.}
\lectureframe{01:09:50}
\end{figure}

\begin{classroomqa}
\audiencequestion{Why is \(\rho^2\), rather than another power, the natural
first correction?}
\lecturerresponse{The ideal term requires one particle to be present in a
small region and therefore scales as \(\rho\).  An interaction requires two
particles to occupy one molecular neighborhood, whose probability scales as
\(\rho^2\).  The factor \(1/2\) is the same pair-counting factor that prevents
\((i,j)\) and \((j,i)\) from being counted separately.
\lecturetimestamp{01:24:27--01:26:00}}
\end{classroomqa}

At fixed temperature one can vary \(V\), measure \(\rho=N/V\), and plot the
pressure.  The ideal law is the linear term.  Curvature of \(P(\rho)\) reveals
the first interaction correction.  Within this weak-potential
approximation, the correction has no explicit temperature dependence because
the factor \(T\) in \eqref{eq:sm6-pressure-from-z} cancels
\(\beta=1/T\).

\begin{classroomqa}
\audiencequestion{How large may the density become before three-body effects
must be included?}
\lecturerresponse{Density counting says that the next virial term is
proportional to \(\rho^3\), but its coefficient cannot be inferred from
\(U_0\) alone.  It receives contributions from correlated products of pair
interactions and, if present, genuine three-body forces.  One must calculate
the next cluster coefficient and compare it with the \(\rho^2\) term; a
guessed formula should not be promoted to a result.
\lecturetimestamp{01:10:50--01:13:38}}
\end{classroomqa}

The ideal-gas term dominates when the magnitude of the correction is small:
\begin{equation}
 \frac{1}{2}\rho^2|U_0|\ll \rho T,
 \qquad\text{or parametrically}\qquad
 |\rho U_0|\ll T.
 \label{eq:sm6-ideal-validity}
\end{equation}
Both sides of the second inequality are energies.  The quantity
\(\rho U_0\) is the interaction-energy scale per particle, while \(T\) is
the kinetic-energy scale.  An ordinary dilute room-temperature gas lies far
inside this regime; substantial deviations arise only when molecular
neighborhoods are occupied often enough for interactions to compete with
thermal motion.

\begin{classroomqa}
\audiencequestion{Can the criterion be estimated for the air in this room?}
\lecturerresponse{Estimate \(U_0\) as a molecular interaction energy times a
molecular volume, multiply by the number density, and compare the result
with \(k_BT\).  The rough molecular numbers show that the interaction scale
is much smaller in ordinary room air.  Considerable compression is required
before molecules frequently crowd into one another's interaction range.
\lecturetimestamp{01:16:24--01:18:54}}
\end{classroomqa}

\begin{classroomqa}
\audiencequestion{May \(U_0\) have either sign?}
\lecturerresponse{Yes.  A positive integrated interaction is net repulsive
in this approximation and raises the pressure above \(\rho T\).  A negative
one is net attractive and lowers it.  This sign test is an immediate physical
check on \eqref{eq:sm6-pressure-correction}.
\lecturetimestamp{01:18:56--01:20:29}}
\end{classroomqa}

\begin{figure}[t]
\centering
\includegraphics[width=0.90\textwidth]{lecture_06_figure_10.png}
\caption{Qualitative repulsive and attractive pair-potential profiles beside
the pressure formula whose correction follows the sign of \(U_0\).}
\lectureframe{01:21:30}
\end{figure}

\begin{classroomqa}
\audiencequestion{Can electrostatic repulsion, gravitational attraction, or
real molecular forces be inserted into this calculation?}
\lecturerresponse{The microscopic potential is input from atomic,
molecular, or quantum physics.  Neutral molecules commonly exhibit a steep
short-range repulsive core and a longer attractive van der Waals tail, often
represented schematically by a 12--6 potential.  The full Mayer-function
formula should be used for such a hard core.  Gravity is different: it is
long-ranged, its spatial integral does not converge in the required way, and
the ordinary short-range thermodynamic limit used here fails.
\lecturetimestamp{01:20:40--01:23:03}}
\end{classroomqa}

The calculation illustrates a useful division of labor.  Microscopic physics
supplies \(u(r)\).  Statistical mechanics turns that energy into \(Z\),
derivatives of \(Z\) produce bulk observables, and physical reasoning checks
their sign, density scaling, extensivity, and domain of validity.  The formal
manipulation is mechanical only after the correct energy and approximation
have been chosen.

\section{Exact Differentials and Path Dependence}

We now change mathematical language.  In the phrase \emph{exact
differential}, exact does not mean numerically precise.  It means that a
one-form is the differential of a genuine scalar function.

Let \(f(x,y)\) be the altitude above a point in a landscape.  A small
displacement changes it by
\begin{equation}
 df=
 \frac{\partial f}{\partial x}\,dx+
 \frac{\partial f}{\partial y}\,dy
 =f_x\,dx+f_y\,dy.
 \label{eq:sm6-total-differential}
\end{equation}
Suppose instead that we are merely handed
\begin{equation}
 \omega=F_x(x,y)\,dx+F_y(x,y)\,dy.
 \label{eq:sm6-one-form}
\end{equation}
Does a scalar \(f\) exist with \(\omega=df\)?  Equality of mixed partials
requires
\begin{equation}
 \frac{\partial F_y}{\partial x}
 =
 \frac{\partial F_x}{\partial y}.
 \label{eq:sm6-exactness-condition}
\end{equation}
On a simply connected domain this zero-curl condition is also sufficient.
The components \((F_x,F_y)\) then form the gradient of \(f\).

\begin{figure}[t]
\centering
\includegraphics[width=0.90\textwidth]{lecture_06_figure_11.png}
\caption{Altitude as a scalar function, its differential, and the
mixed-partial condition for exactness.}
\lectureframe{01:34:30}
\end{figure}

The same condition has an integral meaning.  If \(\omega=df\), then
\begin{equation}
 \int_{\gamma:1\to2}\omega=f(2)-f(1)
 \label{eq:sm6-path-independent-integral}
\end{equation}
depends only on the endpoints, not on the path \(\gamma\).  In particular,
\begin{equation}
 \oint_\gamma\omega=0
 \label{eq:sm6-closed-loop-zero}
\end{equation}
for every contractible closed loop.

Altitude has this property.  Driving between the same two points by a level
road or over a hill gives the same net change in altitude.  Fuel use does
not.  A hilly route and a level route generally consume different amounts,
and a closed trip does not restore the fuel that was burned.  The amount of
fuel remaining is certainly a well-defined quantity in an enlarged state
description, but it is not a function of the car's position \((x,y)\) alone;
it carries memory of the route.  Tiredness supplies the same kind of
analogy: descending a hill does not generally undo the fatigue of climbing
it.

\begin{classroomqa}
\audiencequestion{Is exactness the same idea as a conservative force, and is
it related to curvature?}
\lecturerresponse{A force field with zero curl is locally conservative and
admits a potential, so the language is closely related.  Parallel transport
and curvature also connect closed-loop changes to local geometric data, but
that is a broader subject.  Here only the elementary statement is needed:
zero curl means local path independence, with the usual global condition on
the domain.
\lecturetimestamp{01:39:22--01:41:07}}
\end{classroomqa}

Two examples make the test concrete.  First take
\begin{equation}
 F_x=y,\qquad F_y=x.
 \label{eq:sm6-exact-example}
\end{equation}
Both cross derivatives equal \(1\), and indeed
\begin{equation}
 \omega=y\,dx+x\,dy=d(xy).
 \label{eq:sm6-exact-example-potential}
\end{equation}
Now change one sign:
\begin{equation}
 F_x=y,\qquad F_y=-x.
 \label{eq:sm6-inexact-example}
\end{equation}
Then
\begin{equation}
 \frac{\partial F_x}{\partial y}=1,
 \qquad
 \frac{\partial F_y}{\partial x}=-1,
 \label{eq:sm6-inexact-cross-derivatives}
\end{equation}
so no scalar \(f(x,y)\) has these two components as its gradient.

\begin{figure}[t]
\centering
\includegraphics[width=0.90\textwidth]{lecture_06_figure_12.png}
\caption{An exact example generated by \(f=xy\) and an inexact example whose
cross derivatives are \(1\) and \(-1\).}
\lectureframe{01:44:50}
\end{figure}

\section{Heat and Work Are Not State Functions}

For a simple compressible system with fixed particle number, two independent
equilibrium variables determine the state.  We may use \((T,V)\), \((E,V)\),
or, especially naturally for the fundamental relation, \((S,V)\).  Thus
\begin{equation}
 E=E(S,V),
 \end{equation}
and its exact differential is
\begin{equation}
 dE=
 \left(\frac{\partial E}{\partial S}\right)_V dS+
 \left(\frac{\partial E}{\partial V}\right)_S dV
 =T\,dS-P\,dV.
 \label{eq:sm6-fundamental-relation}
\end{equation}

\begin{classroomqa}
\audiencequestion{Does adiabatic mean slow, insulated, reversible, or
constant entropy?}
\lecturerresponse{The terminology is used differently in different
contexts.  For the equilibrium path considered here, the change is
quasistatic and thermally insulated, hence reversible and isentropic.  A
merely insulated irreversible process can still produce entropy, while the
quantum adiabatic theorem emphasizes slow parameter variation.
\lecturetimestamp{01:47:13--01:47:54}}
\end{classroomqa}

At fixed entropy, an expansion changes the gas energy by
\(-P\,dV\).  At fixed volume, a reversible transfer of heat changes it by
\(T\,dS\).  It is therefore useful to write
\begin{equation}
 \delta W_{\mathrm{on}}=-P\,dV,
 \qquad
 \delta Q_{\mathrm{rev}}=T\,dS,
 \qquad
 dE=\delta Q_{\mathrm{rev}}+\delta W_{\mathrm{on}}.
 \label{eq:sm6-first-law-increments}
\end{equation}
The symbol \(\delta\), rather than \(d\), anticipates that heat and work are
path increments, not differentials of state functions.

\begin{figure}[t]
\centering
\includegraphics[width=0.90\textwidth]{lecture_06_figure_13.png}
\caption{The first law for a piston system:
\(dE=-P\,dV+T\,dS\).}
\lectureframe{01:50:28}
\end{figure}

\begin{classroomqa}
\audiencequestion{Does \(-P\,dV\) represent work done on the gas or by the
gas?}
\lecturerresponse{With the convention in
\eqref{eq:sm6-first-law-increments}, it is work done on the gas.  During
expansion \(dV>0\), the gas does positive work on the piston, so the work done
on the gas is negative and the gas loses energy.  The two descriptions are
equal and opposite.
\lecturetimestamp{01:51:23--01:52:27}}
\end{classroomqa}

The equality \(\delta Q_{\mathrm{rev}}=T\,dS\) must not be extended
unchanged to arbitrary irreversible heating.  In general,
\begin{equation}
 dS=\frac{\delta Q}{T_{\mathrm{boundary}}}+dS_{\mathrm{prod}},
 \qquad dS_{\mathrm{prod}}\ge 0.
 \label{eq:sm6-entropy-production}
\end{equation}
The exact-differential argument below concerns quasistatic equilibrium paths,
for which the reversible relation is the appropriate one.

\begin{classroomqa}
\audiencequestion{Are all paths in the state-space argument assumed to pass
through equilibrium states?}
\lecturerresponse{Yes.  The operations are performed slowly enough that
\(P\), \(T\), \(E\), \(S\), and \(V\) remain well-defined along the path.
The path dependence of heat and work does not require the system to be driven
far from equilibrium.
\lecturetimestamp{01:54:30--01:54:39}}
\end{classroomqa}

To test heat for exactness, use \((E,V)\) as coordinates and rearrange the
first law:
\begin{equation}
 \delta Q_{\mathrm{rev}}=dE+P(E,V)\,dV.
 \label{eq:sm6-heat-one-form}
\end{equation}
If there were a state function \(Q(E,V)\) with \(dQ=\delta Q_{\mathrm{rev}}\),
its partial derivatives would have to be
\begin{equation}
 \left(\frac{\partial Q}{\partial E}\right)_V=1,
 \qquad
 \left(\frac{\partial Q}{\partial V}\right)_E=P.
 \label{eq:sm6-hypothetical-heat-derivatives}
\end{equation}
Equality of mixed derivatives would then require
\begin{equation}
 \frac{\partial}{\partial V}
 \left(\frac{\partial Q}{\partial E}\right)_V
 =
 \frac{\partial}{\partial E}
 \left(\frac{\partial Q}{\partial V}\right)_E,
 \qquad\Longrightarrow\qquad
 0=\left(\frac{\partial P}{\partial E}\right)_V.
 \label{eq:sm6-heat-exactness-requirement}
\end{equation}

\begin{classroomqa}
\audiencequestion{Is this precisely the same curl test used for
\(F_x\,dx+F_y\,dy\)?}
\lecturerresponse{Yes.  The coefficients of \(dE\) and \(dV\) are \(1\) and
\(P(E,V)\).  Their cross derivatives are \(0\) and
\((\partial P/\partial E)_V\).  Heat would be exact only if those derivatives
agreed.
\lecturetimestamp{01:54:39--01:54:52}}
\end{classroomqa}

\begin{figure}[t]
\centering
\includegraphics[width=0.90\textwidth]{lecture_06_figure_14.png}
\caption{The heat one-form \(dE+P\,dV\) and the mixed derivatives required
if it were the differential of a state function.}
\lectureframe{01:58:10}
\end{figure}

The requirement is already physically implausible: adding energy to a gas at
fixed volume generally changes its pressure.  The ideal gas supplies a direct
calculation:
\begin{equation}
 PV=NT,
 \qquad
 E=\frac{3}{2}NT
 \quad\Longrightarrow\quad
 P=\frac{2E}{3V}.
 \label{eq:sm6-ideal-pressure-energy}
\end{equation}
Therefore
\begin{equation}
 \left(\frac{\partial P}{\partial E}\right)_V
 =\frac{2}{3V}\ne0.
 \label{eq:sm6-heat-curl-fails}
\end{equation}
The curl test fails: \(\delta Q_{\mathrm{rev}}\) is not exact.  The work
increment \(\delta W_{\mathrm{on}}=-P\,dV\) fails for the same reason.

\begin{figure}[t]
\centering
\includegraphics[width=0.90\textwidth]{lecture_06_figure_15.png}
\caption{The ideal-gas relations combined into
\(P=2E/(3V)\).}
\lectureframe{02:00:00}
\end{figure}

\begin{figure}[H]
\centering
\includegraphics[width=0.90\textwidth]{lecture_06_figure_16.png}
\caption{The completed conclusion:
\((\partial P/\partial E)_V\ne0\), so the heat one-form is inexact.}
\lectureframe{02:00:35}
\end{figure}

For a closed equilibrium cycle,
\begin{equation}
 \oint dE=0,
 \qquad
 \oint\delta Q_{\mathrm{rev}}
 =-\oint\delta W_{\mathrm{on}}.
 \label{eq:sm6-cycle-balance}
\end{equation}
The system returns to the same energy, but heat and work need not vanish
separately.  Their values depend on the route through state space.  This is
why a cyclic engine can absorb heat and deliver net work even though its
working substance returns to its initial thermodynamic state.

\begin{classroomqa}
\audiencequestion{Is this path dependence connected with engine efficiency,
hysteresis, and memory?}
\lecturerresponse{Engine cycles use it directly: the closed-cycle heat input
balances the net work output while the state function \(E\) returns to its
initial value.  Hysteresis and memory also involve path dependence, although
they often require additional internal variables or nonequilibrium physics.
The present conclusion is narrower and already holds for reversible
equilibrium paths: heat and work are transfers, not properties stored in the
thermodynamic state.
\lecturetimestamp{02:02:21--02:03:17}}
\end{classroomqa}

The two halves of the argument now meet.  Measurement is meaningful only
after its controls and equilibrium conditions are specified.  Approximation
is meaningful only after its small parameter and failure criterion are
specified.  A thermodynamic quantity is a state function only after its
differential passes the path-independence test.  In each case, disciplined
mathematics does not replace physical interpretation; it tells us exactly
where that interpretation is justified.
